\section{The bounded-z/B subdomain of R7 is admitted for all tests}
\label{sec:r17-r7-bounded-a}
\begin{theorem}
Normalize $d=1$ and retain every original R7 condition: $B,z,y>0$
infinitesimal; $t=y/B$ a positive bounded unit; $h=H/B$ and $g=G/B$
infinitesimal; the exclusion of simultaneously infinitesimal $r,j$;
original physicality, stationary contact, coefficient signs, strict gaps,
total-scale intersection and prior exclusions. If $a=z/B$ is bounded,
the entire original parent is finite tensor SOS at every independent
projective testing pair. Both signs, literal zero/infinite angles and
complex states are included. The unbounded-$a$ R7 endpoint remains open.
\end{theorem}

\paragraph{Complete source entry.}
Write $e=1-r$ and $l=1-h$. Since $1-g=a(e-H)$, bounded $a$ is equivalent
on this source shape to $e$ being a positive unit; $a_0=1/e_0$.
We use the letters $e\in(0,1]$ and $t\ge1$ for their real residues.
The strict gap $j=1+zr-l/t$ gives $j_0=1-1/t$. The ordered R7 exclusion
says that $e=1$ and $t=1$ cannot hold together. The original source formulas
give jointly
\[
 C_0={t-e\over2(1+e)}>0,\quad
 (Ah)_0={t\over2(1+e)}>0,\quad
 (\alpha/B^2)_0={t-e\over2e(1+e)}>0,\quad
 \kappa=(mg)_0={1\over2(1+e)}.                        \tag{R7.1}
\]
Thus $A,m$ are unbounded, $C$ and $\alpha/B^2$ are positive units,
and $\alpha$ is infinitesimal. If $e<1$, then
\[
 n_0={e^2\over2(1-e^2)}>0;
\]
if $e=1$, $r$ is infinitesimal and $n$ is unbounded, with $(nr)_0=1/4$.
If $t>1$, then $D_0=1/[2(1+e)(t-1)]$ is a positive unit. If $t=1$,
then $j$ is infinitesimal and $D$ unbounded with $(Dj)_0=\kappa$;
in this case $e<1$ by the retained exclusion. These are exhaustive
source alternatives, with no rate restriction. They also show that
literal $n=0$ is impossible on this admitted subdomain and that total
scale is positive infinitesimal. The $n=0$ boundary on the unbounded-$a$
complement is not removed by this statement.

\paragraph{A complete D-event gauge, including its exact zero syzygy.}
We use the full WA parent, now in the roots
\[
\begin{gathered}
 {\cal R}=(\xi,\eta,R_1,R'_2,T)^T,
 \quad\xi=X+ruZ,\quad\eta=Y-huW,\\
 R_1=uW-vZ/t,\quad R'_2=avW-uZ,\quad T=uvW-BZ.
\end{gathered}
\]
Let $G$ be the exact full-parent matrix obtained from WA8 by leaving
$\eta$ unnormalized and scaling its fourth root by $a$.
The following gauge is the exact D identity developed in the parallel
R8 work; it is valid before any shape specialization. Set
\[
\begin{split}
 d_D&=(uv,-1,l,ruv,-arv)^T,\\
 S_D&=(0,0,(aB-U)v,-(B-V/t)u,U-aV/t)^T,\\
 L_D&=(-2rtu,0,0,-r^2tu,ar^2t)^T.
\end{split}
\]
Then $S_D^T{\cal R}=0$ identically, and
\[
 d_D^T{\cal R}=J+jvZ=:R_{Dc},\qquad
 J=uW+uvX-Y-vZ.                                     \tag{R7.2}
\]
In particular this is the actual contact-zero D row, not an uncentered
positive packet. The exact matrix identity is
\[
 \partial_DG=d_Dd_D^T+\tfrac12(S_DL_D^T+L_DS_D^T)+jH_D,\tag{R7.3}
\]
where the only nonzero entries of $H_D$ are
\[
 (H_D)_{13}=tuv,\quad(H_D)_{15}=-tu,\quad
 (H_D)_{34}=rtuv,\quad(H_D)_{35}=-artv
\]
and their symmetric partners. Consequently the full original parent is
\[
 F/4=D R_{Dc}^2+{\cal R}^T(G|_{D=0}+DjH_D){\cal R}. \tag{R7.4}
\]
The notation $G|_{D=0}$ keeps all other actual coefficients, including
$\alpha/B^2$, unchanged; no modified physical source is asserted.
All full first normals and hard entries are included in R7.3.
This gauge has been independently checked here against the entire WA
matrix, its zero syzygy and the original four-state D event.

\paragraph{One strict local Gram for all source alternatives.}
First apply R7.4, and then retain the centered m excess from R5.2--R5.4,
with positive real cap $K$. The m identity is unaffected by the D gauge:
$H_D$ and its syzygy are independent of $m$ and the D row is independent
of $m$. The remaining full Gram has bounded coefficients after
\[
 \widehat\eta=\sqrt A\,\eta,
\]
and, if $r$ is infinitesimal, also $\widehat\xi=\sqrt n\,\xi$.
Indeed $Ah,nr,mg,Dj$ are bounded, $a,t,l^{-1}$ are bounded units,
$Ah^2$ is infinitesimal, and every large m or D coefficient has been
removed in its complete contact-zero square. No ratio between $A,n,m,D$
is used. This bounds the entire scalar block, both full normal rows and
the hard block, not only the graph.

At $u=v=0$, all entries of $H_D$ vanish. The $\widehat\eta$ pivot is
$1$, its couplings are infinitesimal, the $R_1$ pivot is $C_0+1/2>0$,
and the $T$ pivot is $(\alpha/B^2)_0>0$. If $r$ is a unit, the remaining
$\xi,R'_2$ block is the pure-source rescaling by $\operatorname{diag}(1,a_0^{-1})$
of R6.2. Its determinant is positive because
$n_0+1/2=\kappa/(1-e)$. If $r$ is infinitesimal, normalization by
$\sqrt n$ instead gives separate positive pivots $1$ and
$(K+2\kappa)/a_0^2$; the cross coefficient is infinitesimal.
Thus the complete local remainder has strict residue for every source
alternative. WA10 gives a finite separated tensor SOS for every
independent infinitesimal or zero $u,v$. Restore both retained exact
squares. This proves the full local parent, including its complete
normal/radial jets, without transferring physicality to a capped source.

\paragraph{Remote entry and a finite-dimensional cap lemma.}
For remote charts use the original states and normalize
$\widehat Y=\sqrt A Y$. If $r$ is infinitesimal, normalize also
$\widehat X=\sqrt nX$. First retain the m excess. If $D$ is unbounded,
retain a D excess too, using the exact full-event identity
\[
 DJ^2-(D-K_D)(J+jvZ)^2
 =K_DJ^2-2(D-K_D)jvZJ-(D-K_D)j^2VZ^2.               \tag{R7.5}
\]
All coefficients in the remaining full parent are bounded. The hard
normalized state pivots are positive and decouple in residue; all other
coefficients are bounded and divided by $\sqrt A$ or $\sqrt n$ in
couplings to these pivots. Thus the soft residue has $Y=0$, and also
$X=0$ when $r$ is infinitesimal.

We repeatedly use the following elementary real fact. If a real quadratic
form $B$ is positive definite on the common kernel of real linear rows
$L_1,\ldots,L_k$, then $B+K\sum L_i^2$ is positive definite for sufficiently
large positive real $K$, after any fixed bounded cross terms vanishing on
that common kernel have been included in $B$. This follows by choosing
real coordinates on the kernel and its orthogonal complement and taking
the Schur complement. Here the relevant original event coefficients are
unbounded, so every such real cap leaves positive excesses. The choice
uses only real testing residues, never a comparison of source and testing
rates. The positive hard pivots can then be included in the same strict
full-parent Gram.

\paragraph{Finite remote tests when j is a unit.}
Let $U=u^2,V=v^2$ refer now to the real testing residues, not both zero.
After $Y=0$, the centered m row has residue
$M=vW-X-uZ$. Write $B_{\rm fin}$ for the entire finite parent with its
$m$ packet removed; the bounded cap correction is $2\kappa vWM$ and
vanishes on $M=0$.

If $r$ is a unit, retain $X$ and restrict by $X=vW-uZ$. The two-state
form on $(W,Z)$ has determinant
\[
 { (1+U)(1+V)\over4(1-e)(1+e)^2(t-1)}
       \{(etU-V)^2+(t-e)^2U^2V\}>0.                 \tag{R7.6}
\]
Its $Z=0$ restriction is positive: the positive $n,C$ packets and the
nonnegative original signed core at $Z=0$ force $W=0$ at every remote
pair. Hence it is positive definite, and the cap lemma gives a strict
full residue.

If $r$ is infinitesimal, then $e=1$, $t>1$ and $X$ is another hard state.
For $v\ne0$, restrict by $W=uZ/v$. The remaining scalar is
\[
 {Z^2\over4(t-1)V}\{(tU-V)^2+(t-1)^2U^2V\}>0.      \tag{R7.7}
\]
For $v=0,u\ne0$, $M=0$ forces $Z=0$ and the scalar on $W$ is
$(C_0+D_0+1/2)UW^2>0$. These include the literal first axis and both
angle signs. Again the cap lemma supplies strict residue for the complete
bounded remainder.

\paragraph{Finite remote tests with both m and D unbounded.}
Here $t=1$, $e<1$, $r=1-e$ is a unit and $n$ is finite. With $Y=0$ the
two centered rows have residues
\[
 M=vW-X-uZ,\qquad J=uW+uvX-vZ.
\]
They have rank two at every finite remote pair. Their common kernel is
spanned in $(W,X,Z)$ by
\[
 (v(1+U),\ V-U,\ u(1+V)).                            \tag{R7.8}
\]
On this vector, the complete finite parent with both large packets
removed has value
\[
 { (1+U)(1+V)\over2(1-e)(1+e)}
       \{(eU-V)^2+(1-e)^2U^2V\}>0.                  \tag{R7.9}
\]
The bounded residue corrections are $2\kappa vWM$ and
$-2\kappa vZJ$; both vanish on that common kernel. The cap lemma therefore
gives positive real caps for a strict three-state soft residue; the hard
$\widehat Y$ pivot completes the full parent. The higher source terms
proportional to $mg^2$ and $Dj^2$ have infinitesimal coefficients and are
retained in the bounded remainder, not discarded from the exact identity.

\paragraph{Every reciprocal chart.}
For unbounded $u$, use $p=1/u$ and the entire polynomial parent $p^2F$.
The m row on $Y=0$ is $-Z$. If $D$ is bounded, its kernel $Z=0$ has
complete residue
\[
 n_0(1+V)X^2+C_0(1+V)W^2+(D_0+1/2)(W+vX)^2\succ0, \tag{R7.10}
\]
with $X$ omitted and its normalized hard pivot retained if $n$ is
unbounded. If $D$ is unbounded, its additional row is $W+vX$; on the
common kernel the finite form is
$(1+V)(n_0+C_0V)X^2>0$. In this case $n_0>0$.
The m and D bounded cross corrections contain $p$ and vanish in residue.
The cap lemma applies for every bounded real $v$, including zero.

For unbounded $v$, use $q=1/v$ and $q^2F$. The m row on $Y=0$ is $W$.
If $D$ is bounded, the complete residue on $W=0$ is
\[
 (n_0+1/2)(1+U)X^2+D_0(uX-Z)^2\succ0,               \tag{R7.11}
\]
again with $X$ replaced by its hard normalized pivot when $n$ is
unbounded. If $D$ is unbounded, the additional row is $uX-Z$ and the
common-kernel form is $(n_0+1/2)(1+U)X^2>0$.
The bounded corrections are proportional to $W^2$ and to $Z(uX-Z)$,
so vanish on the relevant common kernels. This covers every bounded
real $u$, including zero.

For both angles unbounded, use $p^2q^2F$ and the exact core R5.16 in
$(W,X,Y,P=pZ,Q=qZ)$, with $Y$ normalized and $X$ normalized if needed.
On the soft states $Y=0$ the m row is $-Q$. For bounded $D$, the finite
residue on $Q=0$ is
\[
 (n_0+1/2)X^2+C_0W^2+D_0(X-P)^2\succ0,              \tag{R7.12}
\]
with the same convention for a hard $X$. If $D$ is unbounded, its row
is $X-P$ and the common-kernel form is
$(n_0+1/2)X^2+C_0W^2>0$. The surviving D correction is proportional to
$P(X-P)$, while the m correction contains $p$; they vanish on the common
kernel. Large real caps give a strict full residue. These identities
are polynomial in $p,q$ and include literal $p=0,q=0$.

\paragraph{Exact join and scope.}
The local, finite remote and three reciprocal configurations exhaust every
independent projective testing pair on the stated source domain. Each
bounded remainder has strict residue and thus the finite separated lift
WA10, while every extracted excess is restored as its exact positive
contact-zero square. The proof uses no mixed inverse or radical; all
normalizations and homogeneous restoration factors are purely source or
purely testing. Every full normal, hard coupling, retained equality and
source opening remains in the identities. Applying the real proof to the
real and imaginary state parts completes the complex assertion. This is
an all-tests admission of bounded-$a$ R7, not an admission of its still
unbounded-$a$ endpoint or of the omitted simultaneous small-$r$, small-$j$
intersection.
